
\documentclass{article}
\usepackage{colortbl}
\usepackage{makecell}
\usepackage{multirow}
\usepackage{supertabular}

\begin{document}

\newcounter{utterance}

\centering \large Interaction Transcript for game `cladder', experiment `full\_v1.5\_default', episode 4700 with Dialogue Architects\_\_qwen3.5{-}9b\_\_SCoRe\_Qwen3.5{-}9B.
\vspace{24pt}

{ \footnotesize  \setcounter{utterance}{1}
\setlength{\tabcolsep}{0pt}
\begin{supertabular}{c@{$\;$}|p{.15\linewidth}@{}p{.15\linewidth}p{.15\linewidth}p{.15\linewidth}p{.15\linewidth}p{.15\linewidth}}
    \# & \multicolumn{2}{c}{Player} && \multicolumn{2}{c}{Game Master} \\
    \hline

    \theutterance \stepcounter{utterance}  
    & & & \multicolumn{4}{p{0.6\linewidth}}{
        \cellcolor[rgb]{0.9,0.9,0.9}{
            \makecell[{{p{\linewidth}}}]{
                \texttt{\tiny{[P1$\langle$GM]}}
                \texttt{You are an expert in causal inference. The following question is not a typical commonsense query, but rather a meticulously designed question created by a professor specializing in causal inference, intended to assess the students' mastery of the course content.} \\
\texttt{Imagine a self{-}contained, hypothetical world with only the following conditions, and without any unmentioned factors or causal relationships: CEO has a direct effect on director and manager. Manager has a direct effect on foot size. Director has a direct effect on foot size. For CEOs who fire employees and managers who don't sign termination letters, the probability of large feet is 62\%. For CEOs who fire employees and managers who sign termination letters, the probability of large feet is 15\%. For CEOs who fire employees and managers who don't sign termination letters, the probability of large feet is 83\%. For CEOs who fire employees and managers who sign termination letters, the probability of large feet is 35\%. The overall probability of CEO's decision to fire the employee is 94\%. Will manager signing the termination letter increase the chance of large feet?} \\
\texttt{Start your answer with 'Yes' or 'No', followed by additional reasoning or evidence to support your explanation.} \\
            }
        }
    }
     \\ \\

    \theutterance \stepcounter{utterance}  
    & \multicolumn{4}{p{0.6\linewidth}}{
        \cellcolor[rgb]{0.561,0.737,0.561}{
            \makecell[{{p{\linewidth}}}]{
                \texttt{\tiny{[P1$\rangle$GM]}}
                \texttt{No} \\
\\ 
\texttt{**Reasoning:**} \\
\\ 
\texttt{To determine if the manager signing the termination letter increases the chance of large feet, we must calculate the causal effect of the manager's action (\$M\$) on the outcome (\$Y\$, large feet). This requires comparing the probability of large feet under the intervention where the manager signs the letter (\$P(Y=1 | do(M=1))\$) versus where they do not (\$P(Y=1 | do(M=0))\$).} \\
\\ 
\texttt{**1. Define the Variables and Structure:**} \\
\texttt{*   \$C\$: CEO fires employee (1 = fires, 0 = does not fire).} \\
\texttt{*   \$M\$: Manager signs termination letter (1 = signs, 0 = does not sign).} \\
\texttt{*   \$Y\$: Large feet (1 = yes, 0 = no).} \\
\texttt{*   **Causal Graph:** \$C \textbackslash{}to M\$, \$C \textbackslash{}to Y\$, \$M \textbackslash{}to Y\$.} \\
\texttt{    *   Note: The prompt states "CEO has a direct effect on director and manager." In the context of the specific probabilities provided regarding "firing employees" and "signing termination letters," the variable \$C\$ (firing) is the common cause influencing both \$M\$ (signing) and \$Y\$ (feet). The "director" mentioned in the setup appears to be a distractor or synonymous with the structural role of the manager in this specific data snippet, as the probabilities are explicitly tied to the Manager's action.} \\
\\ 
\texttt{**2. Analyze the Provided Probabilities:**} \\
\texttt{The prompt provides four conditional probabilities. We must map them to the joint distribution \$P(C, M)\$ to identify which scenarios correspond to which probabilities.} \\
\texttt{*   **Scenario A:** "CEOs who fire (\$C=1\$) and managers who don't sign (\$M=0\$)" \$\textbackslash{}to\$ \$P(Y=1 | C=1, M=0) = 62\textbackslash{}\%\$.} \\
\texttt{*   **Scenario B:** "CEOs who fire (\$C=1\$) and managers who sign (\$M=1\$)" \$\textbackslash{}to\$ \$P(Y=1 | C=1, M=1) = 15\textbackslash{}\%\$.} \\
\texttt{*   **Scenario C:** "CEOs who don't fire (\$C=0\$) and managers who don't sign (\$M=0\$)" \$\textbackslash{}to\$ \$P(Y=1 | C=0, M=0) = 83\textbackslash{}\%\$.} \\
\texttt{*   **Scenario D:** "CEOs who don't fire (\$C=0\$) and managers who sign (\$M=1\$)" \$\textbackslash{}to\$ \$P(Y=1 | C=0, M=1) = 35\textbackslash{}\%\$.} \\
\\ 
\texttt{*Verification of Marginal \$P(C)\$:*} \\
\texttt{The prompt states \$P(C=1) = 94\textbackslash{}\%\$ (Overall probability of firing).} \\
\texttt{Let's check the consistency of the scenarios with this marginal.} \\
\texttt{\$P(C=1) = P(C=1, M=0) + P(C=1, M=1)\$.} \\
\texttt{\$P(C=0) = P(C=0, M=0) + P(C=0, M=1) = 6\textbackslash{}\%\$.} \\
\\ 
\texttt{**3. Calculate the Causal Effect of Manager Signing (\$do(M=1)\$):**} \\
\texttt{We need to calculate \$P(Y=1 | do(M=1))\$ and \$P(Y=1 | do(M=0))\$.} \\
\texttt{Using the backdoor adjustment formula (controlling for the confounder \$C\$):} \\
\texttt{\$\$P(Y=1 | do(M)) = \textbackslash{}sum\_\{c\} P(Y=1 | C=c, M) \textbackslash{}cdot P(C=c)\$\$} \\
\\ 
\texttt{**Calculate \$P(Y=1 | do(M=1))\$ (Manager Signs):**} \\
\texttt{This is the weighted average of the outcome when the manager signs, across both CEO firing statuses.} \\
\texttt{*   If CEO fires (\$C=1\$, prob 0.94): \$P(Y=1 | C=1, M=1) = 0.15\$} \\
\texttt{*   If CEO does not fire (\$C=0\$, prob 0.06): \$P(Y=1 | C=0, M=1) = 0.35\$} \\
\\ 
\texttt{\$\$P(Y=1 | do(M=1)) = (0.15 \textbackslash{}times 0.94) + (0.35 \textbackslash{}times 0.06)\$\$} \\
\texttt{\$\$P(Y=1 | do(M=1)) = 0.141 + 0.021 = 0.162 \textbackslash{}quad (16.2\textbackslash{}\%)\$\$} \\
\\ 
\texttt{**Calculate \$P(Y=1 | do(M=0))\$ (Manager Does Not Sign):**} \\
\texttt{This is the weighted average of the outcome when the manager does not sign.} \\
\texttt{*   If CEO fires (\$C=1\$, prob 0.94): \$P(Y=1 | C=1, M=0) = 0.62\$} \\
\texttt{*   If CEO does not fire (\$C=0\$, prob 0.06): \$P(Y=1 | C=0, M=0) = 0.83\$} \\
\\ 
\texttt{\$\$P(Y=1 | do(M=0)) = (0.62 \textbackslash{}times 0.94) + (0.83 \textbackslash{}times 0.06)\$\$} \\
\texttt{\$\$P(Y=1 | do(M=0)) = 0.5828 + 0.0498 = 0.6326 \textbackslash{}quad (63.26\textbackslash{}\%)\$\$} \\
\\ 
\texttt{**4. Conclusion:**} \\
\texttt{The causal effect of the manager signing the letter is:} \\
\texttt{\$\$Effect = P(Y=1 | do(M=1)) {-} P(Y=1 | do(M=0))\$\$} \\
\texttt{\$\$Effect = 16.2\textbackslash{}\% {-} 63.26\textbackslash{}\% = {-}47.06\textbackslash{}\%\$\$} \\
\\ 
\texttt{The probability of large feet decreases significantly when the manager signs the termination letter. This is because the CEO's decision to fire (which strongly causes large feet in this world) is positively correlated with the manager *not* signing the letter. When the manager signs, it implies the CEO likely did not fire (or the specific causal pathway from firing is absent), leading to a much lower probability of large feet.} \\
\\ 
\texttt{Therefore, manager signing the termination letter **decreases** the chance of large feet.} \\
\\ 
\texttt{**Answer:** No} \\
            }
        }
    }
    & & \\ \\

    \theutterance \stepcounter{utterance}  
    & & & \multicolumn{2}{p{0.3\linewidth}}{
        \cellcolor[rgb]{0.9,0.9,0.9}{
            \makecell[{{p{\linewidth}}}]{
                \texttt{\tiny{[GM$|$GM]}}
                \texttt{No} \\
            }
        }
    }
    & & \\ \\

    \theutterance \stepcounter{utterance}  
    & & & \multicolumn{2}{p{0.3\linewidth}}{
        \cellcolor[rgb]{0.9,0.9,0.9}{
            \makecell[{{p{\linewidth}}}]{
                \texttt{\tiny{[GM$|$GM]}}
                \texttt{no} \\
            }
        }
    }
    & & \\ \\

    \theutterance \stepcounter{utterance}  
    & & & \multicolumn{2}{p{0.3\linewidth}}{
        \cellcolor[rgb]{0.9,0.9,0.9}{
            \makecell[{{p{\linewidth}}}]{
                \texttt{\tiny{[GM$|$GM]}}
                \texttt{game\_result = WIN} \\
            }
        }
    }
    & & \\ \\

\end{supertabular}
}

\end{document}
